\section{Baza danych NoSQL, idea i zastosowania}
Termin NoSQL, często rozwijany jako Not Only SQL, odnosi się do dużej ilości klasy systemów zarządzania bazami danych, które nie korzystają z tradycyjnego modelu relacji między danymi, który dominuje w informatyce od lat siedemdziesiątych XX wieku. Ów model nie powstał w celu wyparcia relacyjnych baz danych, lecz był odpowiedzią na specyficzne wyzwania nowoczesnych aplikacji i systemów Big Data, którym klasyczne bazy danych nie były w stanie sprostać.
\subsection{Geneza: Problem ''One Size Fits All''}
Przez lata relacyjne bazy danych były traktowane jako uniwersalne rozwiązanie (''OneSize Fits All''). Wraz z nadejściem ery Web 2.0, mediów społecznościowych oraz Internetu Rzeczy (IoT) pojawiły się wyzwania określone jako model \textbf{3V}:
\begin{itemize}
    \item \textbf{Objętość (Volume):} Znaczny przyrost danych liczony w petabajtach. W aplikacjach typu \textbf{Multi-Tenant}, czyli aplikacji z jedną instancją i dzierżawą przez wielu klientów, przechowuje się niewyobrażalnie duże ilości danych. Przykładem jest Jira.
    \item \textbf{Szybkość (Velocity):} Potrzeba przetwarzania strumieni danych w czasie rzeczywistym. Na przykład w aplikacjach Internetu Rzeczy, gdzie każde opóźnienie w znaczący sposób wpływa na działanie aplikacji.
    \item \textbf{Różnorodność (Variety):} Aplikację zaczęły mieć potrzebę przechowywania danych w inny sposób niż strukturalny (tabelaryczny). Pojawiły się takie formaty jak JSON i XML, które należą do danych półstrukturalnych jak i dane niestrukturalne jak wideo czy tekst.
\end{itemize}
Fundamentem idei NoSQL jest odejście od sztywnych zasad ACID (Niepodzielności, Spójności, Izolacji i Trwałości) na rzecz wydajnośc i elastyczności. Głównymi założeniami są:
\begin{itemize}
    \item \textbf{Skalowalność horyzontalna:} tradycyjne bazy SQL skalują się głównie wertykalnie, czyli zwiększają swoją wydajność poprzez dokupienie mocniejszego serwera. Takie rozwiązanie jest kosztowne i po siada fizyczne limity. Nierelacyjne bazy danych projektowane są do skalowania horyzontalnego, czyli dołączania kolejnych serwerów do klastra. Takie rozwiązanie jest o wiele tańsze i nie posiada fizicznych limitów.
    \item \textbf{Elastyczność modelu danych:} w relacyjnych systemach zmiana struktury danych, na przykład dodanie nowej kolumny w tabeli, która zawiera miliardy rekordów jest operacją kosztowną i blokującą system. Nierelacyjne bazy danych pozwalają na zapisywanie danych o dowolnej strukturze bez potrzeby definiowania ich schematu. Takie rozwiązanie w znaczący sposób ułatwia tworzenie aplikacji, które posiadają ciągle zmienny format danych.
\end{itemize}
Ze względu na swoje działanie i ograniczenia bazy NoSQL nie są w stanie w pełni zastąpić relacyjnych baz. Systemy bankowe wymagają ścisłych transakcji. Aczkolwiek mogą uzupełnić je w takich obszarach jak:
\begin{itemize}
    \item \textbf{Big Data i Analityka:} Przechowywanie i analiza ogromnych zbiorów danych takich jak logi systemowe czy zachowania użytkowników.
    \item \textbf{Zarządzanie treścią i katalogami produktów:} Elastyczne struktury dokumentowe są idealnym rozwiązaniem dla systemów e-commerce, gdzie każdy produkt może zawierać inny zestaw atrybutów i danych.
    \item \textbf{Internet Rzeczy:} Zbieranie milionów rekordów z czujników na sekundę, gdzie priorytetem jest szybkość zapisu, byłoby niemożliwe przy relacyjnych bazach danych.
    \item \textbf{Media społecznościowe i rekomendacje:} Złożone relacje między użytkownikami, które doprowadziły by do tak zwanej bomby złączeń (JOIN Bomb).
\end{itemize}
Współczesna architektura aplikacja promuje podejście użycia zarówno relacyjnych jak i nierelacyjnych baz danych. Na przykład w nierelacyjnej bazie danych Redis trzyma się sesje użytkowników, w nierelacyjnej dokumentowej bazie MongoDB przechowuje się katalogi produktów, a w relacyjnej bazie PostgreSQL obsługuje się transakcje finansowe.
\section{Klasyfikacja baz NoSQL}
\subsection{Dokumentowe bazy danych}
Głównym założeniem baz dokumentowych jest odejście od struktury tabelarycznej na rzecz elastycznych dokumentów, które pozwalają na przechowywanie danych w formacie zbliżonym do struktur wykorzystywanych w nowoczesnych językach programowania. Dokumenty są pojemnikami dla par klucz-wartość, gdzie wartościami mogą być typy danych pokroju liczby, ciągi znaków, tablice lub zagnieżdżone dokumenty. Większość systemów typu dokument wykorzystuje formaty wywodzące się z języka programowania JavaScript:
\begin{itemize}
    \item \textbf{JSON (JavaScript Object Notation):} Format tekstowy przystępny dla człowieka, który jest standardem stosowanym w komunikacji webowej.
    \item \textbf{BSON (Binary JSON):} Binarny odpowiednik JSON, który jest głównie wykorzystywany między innymi przez MongoDB. Oferuje większą wydajność parsowania oraz rozszerza format JSON oferując dodatkowe typy danych takie jak daty czy dane binarne.
\end{itemize}
Dokumenty o podobnej strukturze są grupowane w kolekcjach, które stanowią odpowiednik tabel w relacyjnych bazach danych. Najważniejszą cechą różniącą nierelacyjne bazy dokumentowe od relacyjnych jest brak konieczności definiowania struktury danych. Każdy dokument moze posiadać inny zestaw par kluczy i wartości. Takie rozwiązanie skutkuje dużą elastycznością i jest idealne w przypadku danych, które ulegają częstym zmianom. Nierelacyjne bazy dokumentowe są głównie projektowanie z myśla o rozproszonej architekturze. Skalowanie odbywa się przede wszystkim na dwóch płaszczyznach:
\begin{itemize}
    \item \textbf{Replikacja:} Tworzenie kopii danych na różnych węzłach. Zapewnia to wysoką dostępność.
    \item \textbf{Sharding:} Horyzontalne partycjoniwanie danych, które polega na rozdzielaniu kolekcji pomiędzy wieloma serwerami na podstawie tak zwanych kluczy shardowania.
\end{itemize}
Nierelacyjne bazy dokumentowe posiadają zarówno zalety jak i wady:
\begin{itemize}
    \item \textbf{Denormalizacja i zagnieżdżenia:} Skutkuje to wysoką wydajnością odczytu ale również redundancja, która zwiększa zapotrzebowanie na miejsce.
    \item \textbf{Elastycznosc schematu:} Pozwala na łatwość iteracji, również otwiera furtkę na niespójność danych.
\end{itemize}
\subsection{Szerokokolumnowe bazy danych}
Szerokokolumnowe bazy danych są zorientowanymi systemami, które charakteryzują się wydajnością przy ogromnych zbiorach danych w architekturach rozproszonych. Pozornie przypominają model relacyjny, ponieważ wykorzystują takie pojęcia jak wiersze i kolumny. Jednakże struktura i sposób składowania danych na dysku jest diamentralnie inna. W przeciwieństwie do systemów zarządzania relacyjną bazą danych, które przechowują dane wierszami, nierelacyjne bazy szerokokolumnowe grupują dane w tak zwane rodziny kolumn.
\begin{itemize}
    \item \textbf{Wiersz:} Każdy wiersz jest identyfikowany przez unikalny klucz. Wiersze w tej samej rodzinie kolumn nie muszą posiadać tych samych zestawów kolumn.
    \item \textbf{Rodzina kolumn:} Logiczna grupa powiązanych ze sobą danych. Są one fizycznie składowane blisko siebie na dysku. Podział na rodziny kolumn ma duże znaczenie dla wydajności operacji zarówno wejścia i wyjścia.
    \item \textbf{Komórka:} Najmniejsza jednostka danych, która jest przecięciem wiersza i kolumny. Posiada mechanizm wersjonowania danych za pomocą znacznika czasu.
\end{itemize}
Większośc baz szerokokolumnowych takich jak Apache Cassandra wykorzystuje strukturę LSM-Tree zamiast tradycyjnych drzew B-Tree.
\begin{itemize}
    \item \textbf{B-Tree} \begin{itemize}
        \item Posiada zbalansowaną drzewiastą strukturę optymalizowaną pod kątem systemów pamięcy masowych.
        \item Wykorzystuje strategię aktualizacji w miejscu. Zmiana rekordu wymaga odnalezienia konkretnej strony na dysku i jej nadpisania.
        \item Minusem tego rozwiązania są wysokie koszty zapisu oraz koniecznością blokowania stron.
        \item Plusem tego rozwiązania jest bardzo szybki odczyt punktowy, co jest znaczącym plusem dla systemów opierajacych się głównie na odczycie.
    \end{itemize}
    \item \textbf{LSM-Tree} \begin{itemize}
        \item Jest zaprojektowana pod kątem maksymalizacji przepustowości zapisu poprzez zmianę losowych operacji wejścia i wyjścia na sekwencyjne.
        \item Dane sa zawsze dopisywane. W tle zachodzi proces scalania plików i usuwania nieaktualnych wersji rekordów.
        \item Największym plusem tego rozwiązania jest ekstremalnie szybki zapis oraz brak fragmentacji przy zapisie. Takie rozwiązanie jest optymalne dla zmiennych dużych zbiorów danych.
        \item Minusem tego rozwiązania jest natomiast nadmiarowość odczytu, ponieważ musi przeszukać parę dysków naraz w celu znalezienia najnowszej wersji. Innym minusem jest również nadmiarowość zapisu, przyczyną takiego zjawiska jest zapisywanie fizycznie na dysku danych wiele razy w celu utrzymania większego porządku, zapisują się małe pliki i łączą w większe, żeby utrzymać porządek.
    \end{itemize} 
\end{itemize}

Proces zapisu przebiega w następujący sposób:
\begin{enumerate}
    \item Nowe dane trafiają do dziennika operacji w celu zapewnienia trwałości i uniknięcia utraty danych.
    \item W kolejnym kroku trafiają do pamięci RAM w odpowiedniej strukturze (\textbf{MemTable}.
    \item Po osiągnięciu określonego rozmiaru w \textbf{MemTable}, dane są zapisywane na dysku jako nieedytowalny plik.
\end{enumerate}
Powyższy proces w bazach szerokokolumnowych oferuje wysoką wydajność zapisu, co czyni je idealnie dostosowanymi do systemów logowania zdarzeń czy do analtyki czasu rzeczywistego.

Takie systemy realizują koncepcję niezależności węzłów w klastrze. Skalowanie odbywanie się horyzontalnie poprzez:
\begin{itemize}
    \item \textbf{Spójne hashowanie:} Pozwala na równomierne rozłożenie danych pomiędzy węzły klastra na podstawie klucza partycji.
    \item \textbf{Replikacja wieloośrodkowa:} Możliwość przechowywania kopii danych w różnych centrach danych, co zapewnia odporność na awarie całych regionów geograficznych.
\end{itemize}
Bazy szerokokolumnowe głównie stosuje się w sytuacjach gdy dane mają charakter szeregów czasowych lub aplikacje potrzebują obsługi miliardów rekordów przy niskich opóźnieniach.


Przykłady baz szerokokolumnowych i ich cechy szczególne:
\begin{itemize}
    \item \textbf{Apache Cassandra:} Architektura pierścienia, która eleminuje pojedyńczy punkt awarii.
    \item \textbf{Google Bigtable:} Usługa chmurowa. Podstawa działania wyszukiwarki Goole i Gmaila.
    \item \textbf{Apache HBase:} Wysoka spójność danych, ścisła integracja z systemem plików HDFS.
    \item \textbf{ScyllaDB:} Wysokowydajny klon Cassandry napisany w języku programowania C++.
\end{itemize}
\subsection{Grafowe bazy danych}
Grafowe bazy danych są zaprojektowne do przechowywania i nawigowania po złożonych relacjach między danymi. W innych modelach NoSQL relacje są symulowane lub pomijane, natomiast w bazach grafowych stanowią one fundament architektury, będąc równie istotne co same węzły.


Większość komercyjnych rozwiązań, na przykład Neo4j, opiera się na modelu etykietowanych właściwości grafu. Składa się z trzech podstawowych elementów:
\begin{itemize}
    \item \textbf{Węzły:} reprezentują encje, na przykład użytkownika, produkt czy miasto. Mogą posiadać etykiety, które kategoryzują ich rolę w grafie.
    \item \textbf{Krawędzie:} definiuja połączenia między węzłami. W bazach grafowych krawędzie zawsze posiadają kierunek jak i posiadają nazwę.
    \item \textbf{Właściwości:} pary klucz-wartość, które mogą być przypisane zarówno do węzłów jak i do krawędzi. Umożliwia to przechowywanie metadanych relacji takich jak datę zawarcia znajomości na krawędzi.
\end{itemize}
Główną różnicą pomiędzy relacyjnymi bazami danych, a natywnymi bazami grafowymi jest jest mechanizm sąsiedztw bezindeksowych. W relacyjnej bazie danych powiązanych rekodów wymagana przeszukania indeksu (złożoność $0(\log n)$, co przy wielokrotnych złączeniach np. \textbf{JOIN} prowadzi do drastycznego spadku wydajności. W natywnej bazie grafowej każdy węzeł przechowuje fizyczne wskaźniki w postaci adresów w pamięci do swoich sąsiadów.

Dzięki sąsiedztw bezindeksowych koszt przejścia grafu zależy wyłącznie od liczby relacji wychodzących z danego węzła, a nie od całkowitego rozmiaru bazy danych. Zapewnia to stałą wydajność ($0(1)$) nawet przy milardach rekordów, co byłoby niemożliwe przy relacyjnych bazach danych.

Ze względu na specyfikę struktury natywnych baz grafowych SQL nie jest efektywny do obsługi grafów. Powstały następujące języki do obsługi zapytań:
\begin{itemize}
    \item \textbf{Cypher:} Język deklaratywyny, który jest inspirowany ASCII-artem. Ma na celu wizualne odwzorowanie ścieżek w grafie. Jest używany w bazie Neo4j.
    \item \textbf{Gremlin:} język imperatywny, służący do programistycznego przechodzenia przez graf. Używany w Apache TinkerPop.
\end{itemize}
Bazy grafowe znajdują zastosowanie w domenach, gdzie analiza powiązań jest ważniejsza niż prosta agregacja danych:
\begin{itemize}
    \item \textbf{Systemy rekomendacyjne:} analiza zakupów w czasie rzeczywistym takich jak "Klienci, którzy kupili to, kupili również te przedmioty..."
    \item \textbf{Wykrywanie oszustw:} śledzenie nietypowych ścieżek przepływu pieniędzy w sieciach bankowych.
    \item \textbf{Zarządzanie tożsamością i dostepem:} modelowanie złożonych struktur organizacyjnych i uprawnień.
    \item \textbf{Sieci społeczne:} Analiza znajomości przyjaciół przyjacieli na przykład osoby, które możesz znać na LinkedIn czy Facebooku.
\end{itemize}
\section{Wady i zalety w porównaniu z bazami relacyjnymi}
Wybór pomiędzy relacyjną a nierelacyjną bazą danych nie jest kwestią wyższości jednego rozwiązania nad drugim, lecz wynika z analizy wymagań projektu.
\subsection{Zalety baz NoSQl względem RDBMS}
\begin{enumerate}
    \item \textbf{Skalowalność horyzontalna:} bazy NoSQL natywnie wspierają partycjonowanie, co pozwala na linoiwy wzrost wydajności poprzez dodawanie tanich serwerów do klastra. Relacyjnie bazy danych głównie skalują się wertykalnie, takie rozwiązanie jest droższe i posiada ograniczenia fizyczne.
    \item \textbf{Elastyczność schematu:} dzięki uproszczonemu modelowi dostępu, na przykład po kluczu głównym i braku narzutu związanego z utrzymywaniem relacji i integralności referencji, nierelacyjne bazy zazwyczaj oferują niższe opóźnienia przy operacjach zapisu i odczytu prostych struktur.
    \item \textbf{Obsługa danych niestrukturyzowanych:} w przypadku danych, które nie pasują do tabel, na przykład dokumenty JSON, grafy powiązań czy strumienie danych z urządzeń IoT aplikacje są zmuszone do rozwiązań nierelacyjnych.
\end{enumerate}

\subsection{Wady i ograniczenia baz NoSQL}
\begin{itemize}
    \item \textbf{Gorsza spójność:} w relacyjnych bazach danych usunięcie rekordu, który jest powiązany z innymi tabelami, jest utrudnione przez klucze obce. Takie systemy wymuszają spójność danych poprzez mechanizmy zabezpieczające, które blokują usunięcie nadrzędnego wpisu, na przykład użytkownika, dopóki istnieją przypisane do niego rekordy, na przykład zamówienia. Wymusza to na programiście ręczne zarządzanie danymi, co w rozbudowanych aplikacjach jest znaczącym utrudnieniem.
    \item \textbf{Brak standaryzacji języka zapytań:} W RDBMS dominuje standard SQL, natomiast w nierelacyjnych systemach każda baza posiada własne API lub język zapytań, co w znaczący sposób zwiększa próg wejścia u utrudnia migrację pomiędzy systemami.
    \item \textbf{Dojrzałośc i ekosystem:} mimo rosnącej popularności wiele rozwiązań NoSQL jest młodszych od relacyjnych odpowiednikó takich jak PostgreSQL czy Oracle, co czasem przekłada się na słabsze narzędzia administracyjne, monitoringowego i mniejszą więdze powszechną.
\end{itemize}
\section{Charakterystyka baz typu key-value}
Systemy bazujące na modelu klucz-wartość stanowią najprostszą, a zarazem fundamentalną formę przechowywania danych w NoSQL.
\subsection{Model logiczny danych}
W przeciwieństwie do relacyjnych baz danych, które narzucają strukturę wierszy i kolumn, bazy klucz-wartość redukują złożonośc modelu danych do zbioru par, gdzie każdemu unikalnemu kluczowi przypisana jest jedna wartość. Formalnie ów model można zapisać jako funkcję mapującą:
\begin{equation}
    f: K \rightarrow V
\end{equation}
gdzie $K$ to zbiór unikalnych identyfikatorów, czyli kluczy, a $V$ to zbiór dowolnych wartości, czyli danych.
\subsection{Klucz i wartość}
Klucz pełni rolę jedynego identyfikatora rekordu. W tym modelu nie występują takie pojecia jak klucz obcy czy relacja między kluczami. Klucz posiada następujące właściwości:
\begin{itemize}
    \item \textbf{Unikalność:} w obrębie jednej przestrzeni nazw klucz musi być unikalny.
    \item \textbf{Rola indeksu:} klucz jest jedynym indeksowanym atrybutem. Skutkuje to tym, że wyszukiwanie rekordów jest możliwe wyłącznie poprzez podanie dokładnego klucza.
    \item \textbf{Partycjonowanie:} klucz determinuje fizyczną lokalizację danych w klastrze.
\end{itemize}
Kluczową cechą wyróżniającą ten model jest traktowanie wartości jako tak zwany nieprzezroczysty blob. Silnik bazy danych nie interpretuje, nie waliduje, ani nie parsuje zawartości wartości. Dla bazy danych wartość jest ciągiem bajtów. Implikacje tego podejścia są następujące:
\begin{itemize}
    \item \textbf{Odpowiedzialność aplikacji:} aplikacja decyduje w jakim formacie są zapisywanie dane. Tak samo jest odpowiedzialna za serializację przy zapisie czy deserializację przy odczycie. Przykładowymi formatami po stronie aplikacji są tekst, JSON, XML, obiekt binarny lub obraz.
    \item \textbf{Brak schematu:} pod różnymi kluczami mogą kryć się dane o zupełnie różnych formatach i rozmiarach.
\end{itemize}
\subsection{Interfejs i operacje}
Ze względu na prostotę modelu, interfejs programistyczny baz klucz-wartość jest minimalistyczny i zazwyczaj ograniczony do trzech podstawowych metod:
\begin{itemize}
    \item \textbf{PUT(klucz, wartość):} zapisanie lub nadpisanie wartości pod danym kluczem.
    \item \textbf{GET(klucz):} pobranie wartości znajdującej się pod kluczem.
    \item \textbf{DELETE(klucz):} usunięcie pary klucz-wartość.
\end{itemize}
Brak języka zapytań jak na przykład SQL eliminuje narzut związany z parsowaniem zapytań i planowaniem ich wykonania co w znaczący sposób wpływa na redukcję opóźnień w systemie. Główną zaletą modelu klucz-wartość jest wydajność. Dzięki zastosowaniu tablic haszujących operacje dostępu do danych charakteryzują się złożonością obliczeniową:
\begin{equation}
    O(1)
\end{equation}
W odróżnieniu od drzew B-Tree stosowanych w bazach relacyjnych (złożoność $0(\log n)$ wydajność baz klucz-wartość nie degraduje się drastycznie wraz ze wzrostem wolumenu danych.
\subsection{Skalowalność i ograniczenia modelu}
Opisywany model jest przystosowani do architektury rozproszonej. Dzięki niezależności par klucz-wartość rekordy mogą myć dzielone pomiędzy wiele węzłów w klastrze. Mechanizm ten zazwyczaj opiera się na spójnym hashowaniu:
\begin{enumerate}
    \item System oblicza skrót z klucza: $H = hash(key)$.
    \item Wartość $H$ determinuje do jakiego węzła trafi dany rekord.
    \item Ponieważ ta operacja jest deterministyczna, aplikacja obsługująca system może samodzielnie obliczyć lokalizację danych ograniczając odpytywanie centralnego serwera zarządzającego.
\end{enumerate}
Optymalizacja pod kątem prostych operacji odczytu i zapisu wiąże się z istotnymi ograniczeniami funkcjonalnymi:
\begin{itemize}
    \item \textbf{Brak zapytań złożonych:} nie jest możliwe wykonywanie zapytań typu zwróć wszystkie rekordy, gdzie wartość > 100, bez pobrania wszystkich kluczy i przetworzenia ich po stronie aplikacji.
    \item \textbf{Brak transakcyjności wielorekordowej:} operacje zazwyczaj są wykonywane w obrębie pojedyńczego klucza. Zmiana wielu kluczy jednocześnie nie gwarantuje spójności ACID. Niektóre zaawansowane systemy implementują mechanizmy transakcyjne, co jest rzadkością w czystym modelu klucz-wartość.
\end{itemize}
\section{Typowe scenariusze użycia baz klucz-wartość}
Ze względu na swoją charakterystykę czyli wysoką wydajność prostych operacji przy jednoczesnym braku zaawansowanych możliwości zapytań, bazy klucz-wartość nie są zazwyczaj stosowane jako główny magazyn danych na systemach produkcyjnych. Pełnia głównie rolę wyspecjalizowanych komponentów rozwiązujących problemy wydajnościowe w architekturze systemów rozproszonych.
\subsection{Pamięć podręczna}
Jest to najbardziej powszechne zastosowanie modelu klucz-wartość. Pamięć podręczna jest wdrażana w celu odciążenia relacyjnych baz danych i przyspieszenia czasu odpowiedzi aplikacji. W tej sytuacji baza klucz-wartość przechowuje wyniki kosztownych operacji na bazie SQL. Proces ten zazwyczaj jest realizowany w następujących krokach:
\begin{enumerate}
    \item Aplikacja otrzymuje żądanie o dane, na przykład statystyki użytkownika sprzed tygodnia.
    \item Sprawdza czy dane istnieją w bazie pod unikalnym kluczem.
    \item Jeżeli dane są dostępne, są zwracane natychmiastowo.
    \item Jeśli danych nie ma pod kluczem, aplikacja pobiera je z głównej bazy SQL, zapisuje je pod kluczem i zwraca klientowi.
\end{enumerate}
Kluczowym mechanizmem wykorzystywanym w wielu przypadkach jest TTL, czyli czas życia klucza i danych. Po upływie danego czasu usuwają się rekordy, co zapobiega przechowywaniu przestarzałych danych.
\subsection{Zarządzanie sesjami użytkowników}
Współczesne aplikacje internetowe, które przechowują sesję używają baz klucz-wartość w celu współdzielenia stanu pomiędzy wieloma serwerami aplikacji. W przypadku mechanizmów równoważących obciążeń zapytań HTTP aplikacja przechowywałaby sesję na paru serwerach, co skutkowałoby brakiem spójności i ciągłym brakiem dostępu do treści w przypadku połączenia ze złym serwerem. Bazy klucz-wartość są idealnym rozwiązaniem dla tego problemu ze względu na:
\begin{itemize}
    \item \textbf{Model dostępu:} sesja jest zawsze identyfikowana przez unikalny klucz, a zawartość to zazwyczaj zserializowany obiekt użytkownika.
    \item \textbf{Niskie opóźnienia:} pobieranie danych sesji musi nastąpić przy każdym połączeniu HTTP, dlatego czas i koszt odczytu jest krytyczny.
\end{itemize}
